\documentclass[11pt]{article}
\usepackage[utf8]{inputenc}
\usepackage[T1]{fontenc}
\usepackage{amsmath,amssymb}
\usepackage[margin=1.1in]{geometry}
\usepackage{authblk}
\usepackage[colorlinks=true,linkcolor=blue,citecolor=blue,urlcolor=blue]{hyperref}

\newcommand{\half}{\tfrac12}
\newcommand{\grad}{\nabla}

\title{Hilbert's Sixth Problem without a Limit:\\
\large Continuum Mechanics from RealQM --- a Programme}

\author[1]{Claes Johnson}
\affil[1]{KTH Royal Institute of Technology, SE-100 44 Stockholm, Sweden \\
\small\texttt{claesjohnson@gmail.com}}

\date{\today}

\begin{document}
\maketitle

\begin{abstract}
Hilbert's sixth problem asks, in the part Hilbert attributed to Boltzmann, for a rigorous passage
from the atomistic description of matter to the equations of continuous media. The classical route
--- Newtonian particles, then Boltzmann's kinetic equation, then Euler and Navier--Stokes --- was
carried through for hard spheres by Deng, Hani and Ma in 2025. That achievement comes with a
structural restriction that is not a defect of the proof but of the route: the Boltzmann--Grad
scaling sends the volume fraction to zero, so what is reached is a \emph{dilute gas}. Liquids,
solids and every dense phase lie outside it by construction. This programme proposes the
complementary route. RealQM has no particles to coarse-grain: it is already a continuum theory, $N$
non-overlapping charge densities on ordinary three-dimensional space, with the interfaces free. So
there is no limit to take, no scaling parameter, and no molecular-chaos hypothesis to justify ---
and correspondingly no dilute restriction. We set out what is already exact, what must be supplied,
and what is genuinely open. Exact and by the chain rule alone: RealQM's localisation energy
\emph{is} the von Weizs\"acker functional, its variation \emph{is} Bohm's quantum potential, and the
theory as used is precisely the zero-velocity equilibrium of a Madelung pair --- continuity plus a
momentum balance whose stress is that potential. To be supplied: the dynamics of that pair per
domain, with the interfaces moving as free surfaces. Genuinely open, and the heart of the
programme: \textbf{dissipation}. Madelung--Euler is non-dissipative, Navier--Stokes is not, and
RealQM has no statistical ensemble from which an $H$-theorem could extract irreversibility. Where
viscosity comes from, if it comes at all, is the question this project exists to answer. We state it
as a question rather than announce an answer.
\end{abstract}

\section{What Hilbert asked, and what has been delivered}\label{sec:ask}

The sixth of Hilbert's 1900 problems is headed \emph{Mathematical Treatment of the Axioms of
Physics}. Its second paragraph is the part at issue here: Hilbert points to Boltzmann's work on the
limiting processes ``which lead from the atomistic view to the laws of motion of continua'' and asks
for them to be made rigorous. The programme it names is therefore specific --- not an axiomatisation
of physics in general, but a derivation of the equations of continuum mechanics from a description
of matter in terms of its constituents.

The classical route runs in two stages. From Newtonian dynamics of many particles to Boltzmann's
kinetic equation, in the Boltzmann--Grad scaling; then from Boltzmann to the fluid equations, in a
hydrodynamic scaling. The first stage is where the difficulty has always been concentrated, because
it must produce irreversible macroscopic behaviour from reversible microscopic laws and must do so
for times long enough to matter. Deng, Hani and Ma completed it for hard spheres on the two- and
three-dimensional torus, obtaining the compressible Euler and incompressible Navier--Stokes--Fourier
equations with quantitative error bounds~\cite{DHM2025}. It is reasonable to call this a solution of
the problem for gas dynamics.

\paragraph{The restriction, which is the opening for this programme.} The Boltzmann--Grad limit takes
$N$ particles of diameter $\varepsilon$ with $N\varepsilon^2$ fixed, so the volume fraction
$N\varepsilon^3 \to 0$. The matter described is therefore arbitrarily dilute, and the critical
literature has made the point sharply: the result characterises a dilute gas and cannot embody the
properties of a dense fluid~\cite{Comment2025}. This is not a gap in the argument to be repaired. It
is what the route is: kinetic theory is a theory of free flight interrupted by rare binary
collisions, and a liquid is not that.

So the honest statement of where the problem stands is narrower than the headline. Continuum
mechanics has been derived from atomism \emph{for gases}. For the condensed phases --- the ones
chemistry, materials science and biology are made of --- it has not, and the classical route does
not reach them.

\section{The inversion}\label{sec:inversion}

RealQM~\cite{realqm,realqmbook} describes an $N$-electron system by $N$ non-negative densities
$\psi_i$, each carrying unit charge on its own domain $\Omega_i$, the domains non-overlapping and
their interfaces free to move to where the energy is least. The energy is Coulomb's:
\begin{equation}\label{eq:E}
E[\{\psi_i\},\{\Omega_i\}]=\sum_i\int_{\Omega_i}\Bigl(\half|\grad\psi_i|^2
-\sum_a\frac{Z_a}{|x-x_a|}\psi_i^2\Bigr)dx
+\half\sum_{i\neq j}\iint\frac{\psi_i^2(x)\psi_j^2(y)}{|x-y|}\,dx\,dy ,
\end{equation}
subject to $\int_{\Omega_i}\psi_i^2=1$.

The observation that organises this programme is that \eqref{eq:E} is not an atomistic description
at all. There are no particles in it. There are continuous fields on ordinary three-dimensional
space, and the only discreteness is in the \emph{number of domains} --- an integer labelling of
regions, not a collection of point masses with positions and momenta.

The consequence for Hilbert's problem is structural. The classical route must build a continuum out
of particles, and every difficulty it faces --- the scaling limit, the propagation of chaos, the
recovery of irreversibility --- arises from that construction. If the starting description is
already a continuum, those difficulties do not arise, because the step they were needed for is not
being taken. What remains is a different and much more tractable question: \emph{are the equations
that \eqref{eq:E} satisfies the equations of continuum mechanics?}

And because no dilution is involved, the question is being asked exactly where the classical route
cannot go. RealQM's established regime is atoms, molecules, condensed phases and
proteins~\cite{realqmbook}. If continuum mechanics follows from \eqref{eq:E}, it follows for dense
matter.

\section{What is already exact}\label{sec:exact}

The first part of the bridge is built and requires no approximation and no limit; it is reported in
full in~\cite{madelung} and summarised here because the programme starts from it.

\paragraph{The localisation energy is the von Weizs\"acker functional.} Writing
$\rho_i=\psi_i^2$ with $\psi_i\ge0$, the chain rule gives $\grad\rho_i = 2\psi_i\grad\psi_i$, hence
\begin{equation}\label{eq:vW}
\half\int|\grad\psi_i|^2\,dx \;=\; \frac18\int\frac{|\grad\rho_i|^2}{\rho_i}\,dx
\;=\; T_{\rm vW}[\rho_i] ,
\end{equation}
in units $\hbar=m=1$. An identity, not a model.

\paragraph{Its variation is the quantum potential.} With $u=\sqrt{\rho}$,
\begin{equation}\label{eq:Q}
\frac{\delta T_{\rm vW}}{\delta\rho} \;=\; -\half\frac{\grad^2\sqrt{\rho}}{\sqrt{\rho}} \;=\; Q ,
\end{equation}
Bohm's quantum potential --- obtained here not by analogy but as the variation of a term
\eqref{eq:E} writes down for an unrelated reason, namely the cost of confining a charge to a region.

\paragraph{The theory as used is a hydrostatics.} Stationarity of \eqref{eq:E} is the statement that
$V+Q$ is constant on each domain, which is the $v=0$ case of the Madelung pair
\begin{align}
\frac{\partial\rho}{\partial t}+\grad\!\cdot(\rho v)&=0 ,\label{eq:cont}\\
\frac{\partial v}{\partial t}+(v\cdot\grad)v&=-\grad\bigl(V+Q\bigr) .\label{eq:euler}
\end{align}
Equation~\eqref{eq:cont} is conservation of charge and \eqref{eq:euler} is a momentum balance whose
stress is supplied by \eqref{eq:Q}. \emph{RealQM is therefore already the equilibrium of a continuum
mechanics, and the constitutive law is not postulated --- it is the variation of the energy the
theory was written with.} This is the sense in which the derivation asked for by Hilbert is, on this
route, partly a matter of recognition rather than of proof.

\paragraph{And the interfaces carry a surface tension.} The free boundary between two domains
contributes an energy per unit area, so the multiphase structure is a continuum mechanics
\emph{with surfaces}, which is what dense matter requires and what a kinetic theory has no room
for~\cite{madelung}.

\section{What must be supplied}\label{sec:supply}

The dynamics. Equations~\eqref{eq:cont}--\eqref{eq:euler} must be posed per domain, with the
interfaces evolving as free surfaces: the kinematic condition that a boundary moves with the normal
velocity of the fluid, and continuity of the chemical potential across it. None of this is new
physics; it is the completion of a pair whose static half is already in use.

Two qualifications belong with it, and they are stated here rather than discovered later.
Madelung hydrodynamics is equivalent to the Schr\"odinger equation only for irrotational $v$ with
quantised circulation, $\oint v\cdot dl = 2\pi n$; both are extra conditions rather than
consequences, and they are where the quantum character re-enters a formulation that otherwise reads
as fluid mechanics. And the free-surface problem for two Madelung fluids in contact is a genuine
free-boundary problem whose well-posedness is not established.

\section{What is genuinely open: dissipation}\label{sec:open}

Here is the hard part, and the reason this is a programme and not a result.

Equations~\eqref{eq:cont}--\eqref{eq:euler} are an \emph{Euler} system. They have no viscosity, no
heat conduction and no entropy production. Navier--Stokes has all three. On the classical route the
irreversibility is manufactured in the first stage: molecular chaos breaks the time symmetry of the
particle dynamics, Boltzmann's $H$-theorem supplies an entropy, and the transport coefficients come
out of the collision operator as integrals over a scattering cross-section. Every one of those
ingredients is statistical, and RealQM has none of them. There is no ensemble, no distribution
function, no collision, and no averaging --- which is precisely why the first stage was free.

The bill arrives here. \textbf{The step from Euler to Navier--Stokes, trivial on the classical route
and the whole content of the transport coefficients, has no mechanism in this framework as it
stands.} Three candidate sources are worth examining, and distinguishing them is the first real task
of the project:

\begin{enumerate}\itemsep3pt
\item \textbf{Interface motion as the dissipative degree of freedom.} A system of $N$ electrons has
$O(N)$ free boundaries. Their motion is not a small correction to the bulk; in this framework
conduction \emph{is} boundary motion~\cite{gas}. Coarse-graining over many interfaces could produce
an effective stress with a dissipative part, in the way that grain-boundary motion produces
macroscopic plasticity. Prediction to test: the resulting viscosity should scale with interface
area per unit volume, hence with $N^{1/3}$ per domain rather than with a mean free path.
\item \textbf{Exchange of territory as an irreversible process.} The partition $\{\Omega_i\}$ is
itself a macroscopic variable. Rearrangements of which electron owns which region are discrete
events with no reverse-path symmetry once the geometry has relaxed, which is a candidate for entropy
production that is geometric rather than statistical.
\item \textbf{Finite precision.} The programme's own thermodynamics locates irreversibility in
finite resolution rather than in statistics~\cite{bodysoul}, and the Clay work treats
well-posedness as a property of computed output~\cite{clay}. If that is the right account, viscosity
in this framework is a residual of unresolved scales and the Euler system is the only exact one ---
a conclusion that would be consistent, but which must be argued, not assumed.
\end{enumerate}

Candidate 3 would make the Euler--Navier--Stokes gap a feature; candidates 1 and 2 would close it.
They make different predictions and are therefore separable, which is the only reason the question
is tractable.

\section{The work programme}\label{sec:programme}

In order, with the first two using machinery that already exists.

\begin{enumerate}\itemsep4pt
\item \textbf{The one-domain dynamics, against a known answer.} Implement
\eqref{eq:cont}--\eqref{eq:euler} on a single domain and check the free-particle case against an
exact Gaussian wave packet. This is the smallest test with a right answer and nothing in the
programme should advance past it until it passes.
\item \textbf{Two domains with a moving interface.} The free-surface problem in its simplest
instance: two densities in contact, interface moving with the normal velocity, chemical potential
continuous. Check against the static answer in the limit of vanishing velocity, which is already
computed.
\item \textbf{The momentum balance as an identity, verified numerically.} Confirm that the stress
in \eqref{eq:euler} integrates to the forces the static solver already reports --- a consistency
check between the hydrodynamic and variational forms of the same theory, and a genuine test because
the two are computed by different code.
\item \textbf{The dissipation question.} Measure whether coarse-graining over interfaces produces a
viscous stress, and whether its magnitude scales with interface area as candidate 1 predicts.
\item \textbf{Equation of state.} Whether the von Weizs\"acker stress gives the right pressure for a
dense phase, which is where this route must beat the dilute-gas result or admit it has not.
\end{enumerate}

\section{What would make this a solution, and what would refute it}\label{sec:settle}

A programme of this kind is worth only as much as its failure conditions, so they are stated in
advance.

It would be \emph{established} if: the one-domain dynamics reproduces free propagation exactly; the
two-domain free-boundary problem is well posed and reduces to the known statics; the momentum
balance agrees with the variational forces; and a dense-phase equation of state comes out within the
model's usual few per cent.

It would be \emph{refuted}, or at least confined, if: the free-boundary problem is ill posed, which
would mean the dynamics cannot be written for more than one electron; or the irrotationality and
circulation-quantisation conditions turn out to be incompatible with the non-overlap constraint,
which would mean the Madelung form and the multiphase form cannot both hold; or the coarse-grained
stress has no dissipative part and finite precision cannot supply one, in which case the programme
reaches Euler and stops --- a real result, but not Navier--Stokes and not a solution of the problem
as posed.

The claim being made at this stage is deliberately modest, and it is a claim about structure rather
than about having done the work. The classical route derives continuum mechanics from atomism for a
dilute gas, by a limit, after a century of effort. This route starts from a description that is
already continuous, so the kinematics and the constitutive law come for nothing, and the entire
difficulty is relocated into a single question --- where dissipation comes from. That is a better
place for the difficulty to be, because it is one question rather than many, and because it can be
answered by computation.

\begin{thebibliography}{9}

\bibitem{DHM2025} Y.~Deng, Z.~Hani and X.~Ma, \emph{Hilbert's sixth problem: derivation of fluid
equations via Boltzmann's kinetic theory}, arXiv:2503.01800 (2025).

\bibitem{Comment2025} \emph{Comment on ``Hilbert's Sixth Problem: Derivation of Fluid Equations via
Boltzmann's Kinetic Theory'' by Deng, Hani and Ma}, arXiv:2504.06297 (2025).

\bibitem{madelung} C.~Johnson, \emph{The Hydrodynamic Form of RealQM: Madelung Hydrostatics, a
Surface Tension, and the Missing Half}.

\bibitem{realqm} C.~Johnson, \emph{RealQM: Quantum Mechanics as Multiphase 3D Continuum Mechanics},
in~\cite{realqmbook}.

\bibitem{realqmbook} C.~Johnson, \emph{Real Quantum Mechanics} (Body\&Soul, Vol.~VI),
\url{https://claes542.github.io/RealMolecule/realquantum3.pdf}.

\bibitem{gas} C.~Johnson, \emph{Conduction in RealQM: the Free Boundary as Charge Carrier},
in~\cite{realqmbook}.

\bibitem{bodysoul} C.~Johnson, \emph{Computational Thermodynamics} (Body\&Soul, Vol.~V).

\bibitem{clay} C.~Johnson, \emph{The Clay Navier--Stokes Problem and Computational
Well-Posedness}.

\end{thebibliography}

\end{document}
